% ======================================================================
%  ÖZEL BÖLÜM: Modern AI Stack — tam sayfa poster
% ======================================================================

\concept{Temel · Foundations}{Modern AI Stack}{Bir AI sisteminin katmanları: tek bir zihinsel modelde}

Bir AI sistemi tek bir ``model çağrısı'' değildir. Bir kullanıcı niyetinden bir
eyleme uzanan bir \emph{katmanlar dizisidir} --- ve her katmanı bu kitapta bir
derin kavram işler. Aşağıdaki poster, kitabın tamamının haritasıdır.

\vspace{4pt}
\begin{center}
\begin{tikzpicture}[font=\sffamily\small,
  layer/.style={draw=accent, line width=0.9pt, fill=wash, rounded corners=3pt,
                minimum width=62mm, minimum height=12mm, align=center},
  flow/.style={-{Stealth[length=2.6mm]}, accent, line width=1pt},
  rail/.style={draw=accentdk, line width=0.8pt, fill=accent!12, rounded corners=3pt,
               minimum width=13mm, minimum height=92mm}]

  % --- merkez akış ---
  \node[layer, fill=white, draw=softgray] (user)
    {\textbf{USER}\\[1pt]{\scriptsize\color{softgray} kullanıcı bir niyetle gelir}};
  \node[layer, below=4mm of user] (intent)
    {\textbf{INTENT}\\[1pt]{\scriptsize\color{softgray} niyeti anla · doğru yola yönlendir}};
  \node[layer, below=4mm of intent] (ctx)
    {\textbf{CONTEXT}\\[1pt]{\scriptsize\color{softgray} doğru bilgiyi pencereye koy}};
  \node[layer, below=4mm of ctx, fill=accentdk, text=white, draw=accentdk] (llm)
    {\textbf{LLM}\\[1pt]{\scriptsize\color{wash} akıl yürüt ve üret}};
  \node[layer, below=4mm of llm] (tools)
    {\textbf{TOOLS}\\[1pt]{\scriptsize\color{softgray} dış dünyada eyle}};
  \node[layer, below=4mm of tools] (mem)
    {\textbf{MEMORY}\\[1pt]{\scriptsize\color{softgray} hatırla · sonraki tura taşı}};
  \node[layer, below=4mm of mem, fill=good!10, draw=good] (out)
    {\textbf{ÇIKTI}\\[1pt]{\scriptsize\color{softgray} cevap veya eylem}};

  \foreach \a/\b in {user/intent, intent/ctx, ctx/llm, llm/tools, tools/mem, mem/out}
    \draw[flow] (\a) -- (\b);

  % --- yan raylar (her katmanı kuşatır) ---
  \node[rail, anchor=north] (ev) at ([xshift=-40mm]user.north)
    {\rotatebox{90}{\textbf{EVALUATION} · kaliteyi ölç}};
  \node[rail, anchor=north] (ob) at ([xshift=40mm]user.north)
    {\rotatebox{90}{\textbf{OBSERVABILITY} · ne olduğunu gör}};

  % --- alt kuşak: guardrails ---
  \node[draw=warn, line width=0.9pt, fill=warmwash, rounded corners=3pt,
        minimum width=136mm, minimum height=8mm, below=5mm of out.south,
        anchor=north, font=\sffamily\small] (g)
    {\textbf{\color{warn}GUARDRAILS \& SAFETY} \,·\, tüm sistemi kuşatır};
\end{tikzpicture}
\end{center}

\vspace{2pt}
{\sffamily\bfseries\small\color{accentdk}Her katman → işleyen derin kavram(lar):}\par
\vspace{2pt}
\noindent\begin{tabularx}{\linewidth}{@{}>{\raggedright\arraybackslash}p{26mm} L@{}}
\arrayrulecolor{hairline}\toprule
\textbf{INTENT} & \term{Prompt} · \term{Router} \\
\textbf{CONTEXT} & \term{Context Engineering} · \term{Embedding} · \term{RAG} · \term{Memory} \\
\textbf{LLM} & \term{Prompt} · \term{Reasoning} · \term{Fine-tuning} \\
\textbf{TOOLS} & \term{Tool Calling} · \term{MCP} \\
\textbf{MEMORY} & \term{Memory} \\
\textbf{Kuşatan} & \term{Eval} · \term{Observability} · \term{Guardrails} \\
\bottomrule
\end{tabularx}

\vspace{4pt}
\begin{principlebox}
\textbf{Stack'i oku:} İstek yukarıdan aşağı akar (kullanıcı → eylem). Eval ve
Observability \emph{her} katmanı dikey olarak kuşatır; Guardrails tüm sistemi
çevreler. Bir AI sorununu teşhis ederken önce ``hangi katman?'' diye sor.
\end{principlebox}
